\subsection{AUTOMORPHISMS OF A CONNECTED COMPACT LIE GROUP}

Denote by \(\mathrm{Aut}(\mathrm{G})\) the Lie group of automorphisms of G (Chap. III, §10, no. 2), and by \(\mathrm{Aut}(\mathrm{G},\mathrm{T})\) the closed subgroup of \(\mathrm{Aut}(\mathrm{G})\) consisting of the elements \(u\) such that \(u(\mathrm{T}) = \mathrm{T}\) . We have seen (§1, no. 4, Cor. 5 of Prop. 4) that the identity component of \(\mathrm{Aut}(\mathrm{G})\) is the subgroup \(\mathrm{Int}(\mathrm{G})\) of inner automorphisms; denote by \(\mathrm{Int}_{\mathrm{G}}(\mathrm{H})\) the image in \(\mathrm{Int}(\mathrm{G})\) of a subgroup H of G.

Let D be the covariant diagram of G relative to T; denote by \(\operatorname{Aut}(D)\) the group of its automorphisms, and by \(\mathrm{W}(D)\) its Weyl group. The map \(u \mapsto \mathrm{D}_{*}(u)\) is a homomorphism from \(\operatorname{Aut}(G, T)\) to \(\operatorname{Aut}(D)\) . Prop. 15 of no. 9 immediately gives:

PROPOSITION 17. The homomorphism \(\operatorname{Aut}(\mathrm{G},\mathrm{T})\to \operatorname{Aut}(\mathrm{D})\) is surjective, with kernel \(\operatorname{Int}_{\mathrm{G}}(\mathrm{T})\) .

Note that \(\operatorname{Aut}(G,T) \cap \operatorname{Int}(G) = \operatorname{Int}_{G}(N_{G}(T))\) and that the image of \(\operatorname{Int}_{G}(N_{G}(T))\) in \(\operatorname{Aut}(D)\) is \(W(D)\) (no. 5, Prop. 10). Thus, Proposition 17 gives an isomorphism

\[
\operatorname{Aut} (\mathrm{G}, \mathrm{T}) / (\operatorname{Aut} (\mathrm{G}, \mathrm{T}) \cap \operatorname{Int} (\mathrm{G})) \rightarrow \operatorname{Aut} (\mathrm{D}) / \mathrm{W} (\mathrm{D}).
\]

Further, \(\operatorname{Aut}(\mathrm{G}) = \operatorname{Int}(\mathrm{G}).\operatorname{Aut}(\mathrm{G}, \mathrm{T})\) . Indeed, if u belongs to \(\operatorname{Aut}(\mathrm{G})\) , \(u(\mathrm{T})\) is a maximal torus of T, hence is conjugate to T, and there exists an inner automorphism v of G such that \(u(\mathrm{T}) = v(\mathrm{T})\) , in other words \(v^{-1}u \in \operatorname{Aut}(\mathrm{G}, \mathrm{T})\) . It follows that \(\operatorname{Aut}(\mathrm{G})/\operatorname{Int}(\mathrm{G})\) can be identified with \(\operatorname{Aut}(\mathrm{G}, \mathrm{T})/(\operatorname{Aut}(\mathrm{G}, \mathrm{T}) \cap \operatorname{Int}(\mathrm{G}))\) , so in view of the preceding we have an exact sequence

\[
1 \to \operatorname{Int} (\mathrm{G}) \to \operatorname{Aut} (\mathrm{G}) \to \operatorname{Aut} (\mathrm{D}) / \mathrm{W} (\mathrm{D}) \to 1.\tag{16}
\]

Consequently:

PROPOSITION 18. The group \(\operatorname{Aut}(\mathrm{G}) / \operatorname{Int}(\mathrm{G})\) is isomorphic to \(\operatorname{Aut}(\mathrm{D}) / \operatorname{W}(\mathrm{D})\) .

In particular, assume that G is semi-simple; the group \(\operatorname{Aut}(D)\) can then be identified with the subgroup of \(\mathrm{A}(\mathrm{R}(\mathrm{G},\mathrm{T}))\) (Chap. VI, §1, no. 1) consisting of the elements u such that \(u(\mathrm{X}(\mathrm{T})) \subset \mathrm{X}(\mathrm{T})\) , and the subgroup \(\mathrm{W}(\mathrm{D})\) can be identified with \(\mathrm{W}(\mathrm{R}(\mathrm{G},\mathrm{T}))\) .

COROLLARY. If G is simply-connected, or if C(G) reduces to the identity element, the group \(\operatorname{Aut}(\mathrm{G})/\operatorname{Int}(\mathrm{G})\) is isomorphic to the group of automorphisms of the Dynkin graph of R(G,T).

This follows from the preceding and Chap. VI, §4, no. 2, Cor. of Prop. 1.

We are now going to show that the extension (16) admits sections.

For all \(\alpha\in\mathrm{R}(\mathrm{G},\mathrm{T})\) , denote by \(\mathrm{V}(\alpha)\) the 2-dimensional vector subspace of g such that \(\mathrm{V}(\alpha)_{(\mathbf{C})}=\mathfrak{g}^{\alpha}+\mathfrak{g}^{-\alpha}\) ; denote by K the quadratic form associated to the Killing form of g.

DEFINITION 3. A framing of (G,T) is a pair \((\mathrm{B},(U_{\alpha})_{\alpha \in \mathrm{B}})\) , where B is a basis of \(\mathrm{R}(\mathrm{G},\mathrm{T})\) (Chap. VI, §1, no. 5, Def. 2) and where, for all \(\alpha \in \mathrm{B}\) , \(U_{\alpha}\) is an element of \(\mathrm{V}(\alpha)\) such that \(\mathrm{K}(U_{\alpha}) = -1\) .

A framing of G is a maximal torus T of G together with a framing of \((G, T)\) .

Lemma 3. Let \(B_{0}\) be a basis of \(\mathrm{R}(\mathrm{G},\mathrm{T})\) . The group \(\operatorname{Int}_{\mathrm{G}}(\mathrm{T})\) operates simply-transitively on the set of framings of \((\mathrm{G},\mathrm{T})\) of the form \((\mathrm{B}_{0},(U_{\alpha})_{\alpha\in\mathrm{B}_{0}})\) .

For all \(\alpha\in B_{0}\) , denote by \(\mathrm{K}(\alpha)\) the restriction of the quadratic form K to \(\mathrm{V}(\alpha)\) ; the operation of T on \(\mathrm{V}(\alpha)\) defines a morphism \(\iota_{\alpha}:T\to\mathbf{SO}(\mathrm{K}(\alpha))\) . We have seen in no. 4 that \(\mathbf{SO}(\mathrm{K}(\alpha))\) can be identified with U in such a way that \(\iota_{\alpha}\) is identified with the root \(\alpha\) . Since \(B_{0}\) is a basis of R, it is a basis of the Z-module \(\mathrm{Q}(\mathrm{R})\) generated by the roots, hence a basis of the submodule \(\mathrm{X}(\mathrm{T}/\mathrm{C}(\mathrm{G}))\) of \(\mathrm{X}(\mathrm{T})\) . It follows that the product of the morphisms \(\iota_{\alpha}\) induces an isomorphism from \(\mathrm{T}/\mathrm{C}(\mathrm{G})\) to the product of the groups \(\mathbf{SO}(\mathrm{K}(\alpha))\) . But the latter group operates simply-transitively on the set of framings of (G,T) whose first component is \(B_{0}\) .

PROPOSITION 19. The group \(\operatorname{Int}(\mathbf{G})\) operates simply-transitively on the set of framings of \(\mathbf{G}\) .

Let \(e = (\mathrm{T}, \mathrm{B}, (U_{\alpha}))\) and \(e' = (\mathrm{T}', \mathrm{B}', (U_{\alpha}')\) be two framings of \(G\) . There exist elements \(g\) in \(G\) such that \((\operatorname{Int} g)(\mathrm{T}) = \mathrm{T}'\) , and these elements form a single coset modulo \(N_G(T)\) . Thus, we can assume that \(T = T'\) , and we must prove that there exists a unique element of \(\operatorname{Int}_G(N_G(T))\) that transforms \(e\) to \(e'\) . By Chap. VI, §1, no. 5, Remark 4, there exists a unique element \(w\) of \(W(R)\) such that \(w(B) = B'\) . Since \(W(R)\) can be identified with \(N_G(T)/T\) , there exists \(n \in N_G(T)\) such that \(w = \operatorname{Int} n\) , and \(n\) is uniquely determined modulo \(T\) . Thus, we can assume that \(B = B'\) , and we must prove that there exists a unique element of \(\operatorname{Int}_{\mathrm{G}}(\mathrm{T})\) that transforms e to \(e'\) , which is simply Lemma 3.

COROLLARY. Let e be a framing of (G,T) and let E be the group of automorphisms of G that leave e stable. Then \(\operatorname{Aut}(G)\) is the semi-direct product of E by \(\operatorname{Int}(G)\) , and \(\operatorname{Aut}(G,T)\) is the semi-direct product of E by \(\operatorname{Int}(G)\cap\operatorname{Aut}(G,T)=\operatorname{Int}_{G}(\mathrm{N}_{\mathrm{G}}(\mathrm{T}))\) .

Indeed, every element of \(\operatorname{Aut}(G)\) transforms e into a framing of G. By Prop. 19, every coset of \(\operatorname{Aut}(G)\) modulo \(\operatorname{Int}(G)\) meets E in a single point, hence the first assertion. The second is proved in the same way.

Remark. Let G and \(G'\) be two connected compact Lie groups, and let \(e = (\mathrm{T}, \mathrm{B}, (U_{\alpha}))\) and \(e' = (\mathrm{T}', \mathrm{B}', (U_{\alpha}')\) be framings of G and \(G'\) , respectively. Let X be the set of isomorphisms from G to \(G'\) that take e to \(e'\) . The map \(f \mapsto \mathrm{D}^{*}(f)\) (resp. \(\mathrm{D}_{*}(f)\) ) is a bijection from X to the set of isomorphisms from \(\mathrm{D}^{*}(\mathrm{G}', \mathrm{T}')\) to \(\mathrm{D}^{*}(\mathrm{G}, \mathrm{T})\) (resp. \(\mathrm{D}_{*}(\mathrm{G}, \mathrm{T})\) to \(\mathrm{D}_{*}(\mathrm{G}', \mathrm{T}')\) ) that map \(B'\) to B (resp. B to \(B'\) ). Indeed, this follows immediately from Prop. 15 and Lemma 3.

